\section{The two pure conductor-sixteen characters at the quarter residues}
\label{sec:pure-conductor16-quartic-strict}

\paragraph{Scope.}
The original finite characters in this argument are the two primitive
odd quartic characters modulo sixteen, with trivial cubic part.
The real quartic class inputs are already proved in AC-G147-class.
Their two imaginary companions have explicit integral Stickelberger
annihilators that are units modulo $269$ on the required sectors.
Two existing global units then give the required local injections.
No new analytic interval or complete unit basis is required.

Put $p=269$, $K=\mathbb Q(\sqrt{-3})$, and fix
$i\in\mathbb Q_p$ with $i^2=-1$ and $i\equiv82\pmod p$.
Let $\phi$ be the rational primitive odd quartic character modulo
sixteen with $\phi(13)=i$ and $\phi(5)=-i$.
The same notation denotes its restriction to $G_K$ when needed.
Write $c_p$ for the cyclotomic character and $\omega_p$ for its
Teichmuller part.  For a discrete coefficient $D$ define
\[
 \begin{gathered}
 S_{\rm str}(K,D)=\ker\bigl(H^1_{\rm cts}(G_K,D)
                                  \longrightarrow\mathcal L(D)\bigr),\\
 \mathcal L(D)=\prod_{v\mid p}H^1(G_{K_v},D)
              \times\prod_{v\nmid p,\infty}H^1(I_v,D).
 \end{gathered}
 \tag{P16Q.1}
\]

\begin{theorem}[AC-G166: the pure conductor-sixteen quarter-residue rows]
\label{AC-G166}
Let $E/\mathbb Q_p$ be any finite extension with its original
ring of integers $O$.  Let $\sigma\in\{1,-1\}$,
$m\in\mathbb Z$, $e\geq1$, and let
$\kappa:G_K\to1+\mathfrak m_O$ be a continuous character
unramified at every finite place away from $p$.
If $e+m\equiv66$ or $200\pmod{268}$, then
\[
 \begin{gathered}
 D=(E/O)(\phi^\sigma\omega_p^{-m}\kappa^{-1}c_p^{-e}),\\
                         S_{\rm str}(K,D)=0.
 \end{gathered}
 \tag{P16Q.2}
\]
The corresponding residual strict groups vanish over every finite
field of characteristic $p$.  The original coefficient ring may
have arbitrary ramification.  For $m=0,\kappa=1$, these are two
actual finite labels in every prescribed host containing
$\mathbb Q(\zeta_{16})$, with no cubic factor.
\end{theorem}

\begin{proof}
\proofstep{1. The two exact fields and the inverse-character dictionary.}
Put $\psi=\omega_p^{67}$, so $\psi(2)=-i$, and set
\[
 \xi_1=\phi\psi,\qquad \xi_2=\phi\psi^{-1},\qquad
 \xi_1^2=\xi_2^2=\chi_8\tau_p.
 \tag{P16Q.3}
\]
These are the two distinct primitive even quartic characters of
conductor $16p$ identified in AC-G147-class.
With $d=\sqrt{2p}$ their real fields and automorphisms are
\[
 \begin{gathered}
 F_i=\mathbb Q(\beta_i),\qquad
 \beta_1^2=2p+23d,\qquad\beta_2^2=2p+3d,\\
 \beta_i^4-4p\beta_i^2+C_i p=0,
       \qquad(C_1,C_2)=(18,1058),\\
 b_i(\beta_i)=\frac{\beta_i^3-A_i\beta_i}{69},
       \qquad(A_1,A_2)=(1067,547),\\
 b_i^2(\beta_i)=-\beta_i,\qquad\xi_i(b_i)=i.
 \end{gathered}
 \tag{P16Q.4}
\]
The actual cyclotomic character construction and these orientations
are those of (G147C.1)--(G147C.8), including the two different kernels.
Both unique quadratic subfields are $Q=\mathbb Q(d)$.
Since the $F_i$ are real, they are disjoint from $K$.
Thus $N_i=KF_i$ has degree eight over $\mathbb Q$ and degree
four over $K$.  Extend $b_i$ to fix $K$, and let $c$ fix $F_i$
and conjugate $K$.  The commuting automorphisms have orders
four and two.  The CRT representative $3229$ is $13$ modulo
sixteen, $1$ modulo three, and $1$ modulo $p$.
It fixes $K$ and induces the stated $b_i$, with character value $i$.

Let $a=\bar\phi^\sigma\bar c_p^{-(e+m)}$ be the residual
coefficient character and put $r\equiv e+m\pmod{268}$.
The Kummer quotient $\eta=a/\bar c_p$ and its inverse have
the following exact dictionary:
\[
 \begin{array}{c|c|c|c|c}
 \sigma&r&\eta&\text{splitting cover}&\eta^{-1}(b_i)\\ \hline
 1&66&\bar\xi_2&N_2/K&187=-\bar i\\
 -1&66&\bar\xi_1^{-1}&N_1/K&82=\bar i\\
 1&200&\bar\xi_1&N_1/K&187=-\bar i\\
 -1&200&\bar\xi_2^{-1}&N_2/K&82=\bar i.
 \end{array}
 \tag{P16Q.5}
\]
Here $-(66+1)=-67$ and $-(200+1)\equiv67\pmod{268}$.
The factors involving $\psi$ occur only in the residual splitting
fields.  The original finite labels remain $\phi$ and $\phi^{-1}$.

\proofstep{2. The two companion class certificates.}
Let $L_i=F_i(\sqrt[3]p)$.  AC-G147-class proves
$\operatorname{Cl}(L_i)[p]=0$.
Norm after extension of ideals is multiplication by three from
$F_i$, and by six from $Q$; hence
\[
                 \operatorname{Cl}(F_i)[p]
                    =\operatorname{Cl}(Q)[p]=0.
 \tag{P16Q.6}
\]
Write $\eta_K=\chi_{-3}$ and $\xi'_i=\eta_K\xi_i$.
These characters are odd and primitive of conductor
$f=48p=12912$: the primitive conductor-three, conductor-sixteen,
and conductor-$p$ factors have coprime moduli.
With $z_i=\sqrt{-3}\,\beta_i$, their fields are
\[
 \begin{gathered}
 F'_i=\mathbb Q(z_i)=N_i^{\langle cb_i^2\rangle},\qquad
 z_i^4+12pz_i^2+9C_i p=0,\\
 b_i(z_i)=-\frac{z_i^3+3A_i z_i}{207},
                        \qquad\xi'_i(b_i)=i.
 \end{gathered}
 \tag{P16Q.7}
\]
The polynomials are Eisenstein at $p$, so they have degree four.
The involution $cb_i^2$ fixes $z_i$ and is precisely the kernel
of the twisted character on $C_2\times C_4$.
Both companion quadratic subfields are again $Q$.

Put $C=\mathbb Q(\zeta_f)$.
Its degree is $4288$, and $[C:F'_i]=1072$; both are prime to $p$.
For $(a,f)=1$ let $\sigma_a(\zeta_f)=\zeta_f^a$.
Use the exact integral annihilator
\[
 \begin{aligned}
 \Theta&=\frac1f\sum_{\substack{1\leq a\leq f\\(a,f)=1}}
                            a\,\sigma_a^{-1},\\
 A&=(5-\sigma_5)\Theta
   =\sum_{\substack{1\leq a\leq f\\(a,f)=1}}
                  \left\lfloor\frac{5a}{f}\right\rfloor\sigma_a^{-1}.
 \end{aligned}
 \tag{P16Q.8}
\]
Reindexing multiplication by five modulo $f$ proves the second
equality over the integers, before any reduction modulo $p$.
Sharifi's Definition~4.5.1, Lemma~4.5.2 and Theorem~4.5.6 give
that $\mathbb Z[\operatorname{Gal}(C/\mathbb Q)]\Theta
\cap\mathbb Z[\operatorname{Gal}(C/\mathbb Q)]$ annihilates
$\operatorname{Cl}(C)$.  The source uses $\sigma_5-5$;
its negative in (P16Q.8) belongs to the same ideal.

The following are exact integer sums, not numerical special-value
approximations.  The finite verifier uses $2$ as a primitive root
modulo $p$, $\psi(2)=-i$, $\eta_K(a)=1,-1$ according as
$a\equiv1,2\pmod3$, and the table
\[
 \begin{array}{c|rrrrrrrr}
 a\bmod16&1&3&5&7&9&11&13&15\\ \hline
 \phi(a)&1&-i&-i&1&-1&i&i&-1.
 \end{array}
 \tag{P16Q.9}
\]
Characters vanish on nonunits.  Independent weighted and floor
summations give
\[
 \begin{gathered}
 \begin{array}{c|c|c}
 i&\sum_{a=1}^f a\xi'_i(a)&B_1(\xi'_i)\\ \hline
 1&464832-154944i&36-12i\\
 2&-309888-258240i&-24-20i
 \end{array}\\[3pt]
 \sum_{(a,f)=1}\left\lfloor5a/f\right\rfloor\xi'_i(a)^{-1}
 =\begin{cases}192+24i,&i=1,\\-100+124i,&i=2.\end{cases}
 \end{gathered}
 \tag{P16Q.10}
\]
Here $B_1(\xi')=f^{-1}\sum a\xi'(a)$, and
$\xi'_i(5)=i$ for both rows.
Thus the displayed floor sum also equals
$(5-i)B_1((\xi'_i)^{-1})$, with the inverse orientation required
by $\sigma_a^{-1}$.
On the two conjugate quartic character components, the integral
annihilator has residues
\[
 \begin{array}{c|cc}
 i&\xi'_i&(\xi'_i)^{-1}\\ \hline
 1&8&107\\
 2&115&223.
 \end{array}
 \tag{P16Q.11}
\]
All four are units modulo $p$.  This calculation does not divide
by the $p$-divisible conductor in a residue field.

Consequently the corresponding components of
$\operatorname{Cl}(C)[p]$ vanish.
Extension from $F'_i$ to $C$ is equivariant and injective on
$p$-torsion, because norm after extension is multiplication
by $1072$.  This kills the $b_i^2=-1$ part of
$\operatorname{Cl}(F'_i)[p]$.
On the $b_i^2=1$ part, norm to $Q$ followed by extension is
$1+b_i^2=2$; its target is zero by (P16Q.6).
Therefore
\[
                         \operatorname{Cl}(F'_i)[p]=0.
 \tag{P16Q.12}
\]
No exact class number, complete unit basis or analytic regulator
calculation is asserted for these companions.

\proofstep{3. The precise class sector and full-local Kummer kernel.}
On $\operatorname{Cl}(N_i)[p]$, the norm-extension operators
to $F_i$ and $F'_i$ are $1+c$ and $1+cb_i^2$.
Both vanish by (P16Q.6) and (P16Q.12).
On $b_i^2=-1$ their sum is multiplication by two, whence
\[
                  \operatorname{Cl}(N_i)[p]^{b_i^2=-1}=0.
 \tag{P16Q.13}
\]
No vanishing of the other class sector of $N_i$ is required.

Let $k$ be a finite coefficient field of characteristic $p$.
Since $[N_i:K]=4$ is prime to $p$, global and local
restriction--inflation identify the residual strict group with
\[
 \begin{gathered}
 e_{\eta^{-1}}\ker\left(V_{N_i,S}\otimes_{\mathbb F_p}k
                       \longrightarrow\mathcal P_{N_i}\otimes_{\mathbb F_p}k\right),\\
 V_{N_i,S}=\{x\in N_i^\times:p\mid v(x)\text{ for }v\nmid p\}/N_i^{\times p},\\
 \mathcal P_{N_i}=\bigoplus_{v\mid p}N_{i,v}^\times/N_{i,v}^{\times p}.
 \end{gathered}
 \tag{P16Q.14}
\]
Indeed the coefficient becomes $k(1)$ over $N_i$; the remaining
quotient twist is $\eta$, so invariants select $\eta^{-1}$.
Every local decomposition and inertia quotient also has order
dividing four.  The condition at $p$ is the full local power
condition, rather than a valuation-only condition.
It makes valuations above $p$ divisible by $p$ as well.
Thus the kernel lies in the all-primes valuation group $V_{N_i}$,
which fits into
\[
 0\longrightarrow\mathcal O_{N_i}^\times/\mathcal O_{N_i}^{\times p}
 \longrightarrow V_{N_i}\longrightarrow\operatorname{Cl}(N_i)[p]
 \longrightarrow0.
 \tag{P16Q.15}
\]
Every character in (P16Q.5) has $b_i^2=-1$.
Equation (P16Q.13) therefore leaves only its unit component.
No adjoining of $\mu_p$ has been used.

\proofstep{4. Two actual global units.}
The original AC-G147 files contain the following units, with no
$t$-coordinates: source unit two for field one, and source unit
three for field two.
\[
 \begin{aligned}
 69U_1={}&-1425171+49828\beta_1+318141\beta_1^2-9725\beta_1^3,\\
 69U_2={}&4653963451966989+188807670726232\beta_2\\
        &-9935544878139\beta_2^2-403077313105\beta_2^3.
 \end{aligned}
 \tag{P16Q.16}
\]
Exact quartic regular-representation matrices give
\[
 \begin{aligned}
 U_1^4-9839692U_1^3+3565381946U_1^2+9839692U_1+1&=0,\\
 U_2^4+40078822751988U_2^3\qquad&\\
 {}-48336775602484788614U_2^2-40078822751988U_2+1&=0.
 \end{aligned}
 \tag{P16Q.17}
\]
For each unit the verifier checks its exact rational matrix
characteristic polynomial and a separate Horner evaluation.
The cube of the displayed polynomial is exactly its original
degree-twelve relation in the unchanged source file.
Monicity and constant coefficient one make each value and its
inverse integral, so these are actual global units independently
of any proposed full unit basis.

\proofstep{5. The full local projectors, with both orientations retained.}
Each polynomial in (P16Q.4) is Eisenstein at $p$.
The completion $K_p$ is the unramified quadratic field.
Hence $N_i$ has one prime over $p$, with ramification degree
four and residue degree two; $\beta_i$ is a uniformizer and
$v_{\beta_i}(p)=4$.
Its power basis over $\mathcal O_{K_p}$ is integral and complete.
All denominators in (P16Q.16) are retained before reduction.
In $\mathbb F_{p^2}[\beta_i]/(\beta_i^4)$ the exact reductions are
\[
 \begin{aligned}
 U_1&\equiv187+36\beta_1+143\beta_1^2+15\beta_1^3,\\
 U_2&\equiv187+212\beta_2+54\beta_2^2+38\beta_2^3,\\
 \ell_1&\equiv262\beta_1+256\beta_1^3,\\
 \ell_2&\equiv168\beta_2+207\beta_2^3,\\
 b_1(\beta_1)&\equiv82\beta_1+39\beta_1^3,\\
 b_2(\beta_2)&\equiv187\beta_2+39\beta_2^3.
 \end{aligned}
 \tag{P16Q.18}
\]
Here $\ell_i=\log(U_i/\operatorname{Teich}(\bar U_i))$.
The local leading eigenvalues differ, although
$\xi_i(b_i)=i$ in both global character fields.
The logarithm is computed by solving $U_i\ell_i'=U_i'$ and
integrating through degree three; all denominators are prime
to $p$.  The nonlinear cubic term of each $b_i$ is retained.

For $\lambda=82,187$, let
$P_\lambda=\sum_{j=0}^3\lambda^{-j}b_i^j$ on unit classes modulo $p$.
Full substitution gives
\[
 \begin{array}{c|c|c}
 i&P_{82}\ell_i&P_{187}\ell_i\\ \hline
 1&241\beta_1+118\beta_1^3&99\beta_1^3\\
 2&148\beta_2^3&134\beta_2+142\beta_2^3.
 \end{array}
 \tag{P16Q.19}
\]
All four values are nonzero and satisfy their stated eigenvalue
identities under the complete local action.
Integer representatives of the projector coefficients define
products of actual global conjugate units before localization.
Changing representatives by multiples of $p$ does not change
the Kummer class or the displayed truncated logarithm.
Since $4<p-1$, logarithm and exponential identify principal
units with the maximal ideal.  A principal $p$th power has
logarithm of valuation at least five, whereas the four images
have leading degree one or three.
Teichmuller units are $p$th powers; a $p$th root of a principal
unit has residue one by bijectivity of Frobenius on the residue
field.  Thus the images are nonzero in the full local Kummer
quotient, not merely in a selected valuation quotient.

\proofstep{6. Unit multiplicity and residual strict zero.}
The four archimedean places above a fixed complex embedding of
$K$ form a regular $C_4=\operatorname{Gal}(N_i/K)$ set.
The equivariant unit character is the regular character minus
the trivial character.  There is no $p$-torsion among the roots
of unity, since $p-1=268$ does not divide $[N_i:\mathbb Q]=8$.
Because $p\nmid4$, each of the eigenvalues $82$ and $187$ has
a one-dimensional unit space over $\mathbb F_p$.
The four nonzero images in (P16Q.19) make localization injective
on every required space, and this remains so after tensoring
with any finite $k$.
Together with (P16Q.13)--(P16Q.15) and the inverse-character
table (P16Q.5), this proves every claimed residual strict zero.

\proofstep{7. The original, arbitrarily ramified coefficient ring.}
Let $\varpi$ be a uniformizer of the original $O$.
The coefficient $D$ in (P16Q.2) is divisible and
$D[\varpi]\simeq k(a)$.
At $p$, the conductor-sixteen finite character is unramified
and $\kappa$ reduces to one.  Since $K_p/\mathbb Q_p$ is
unramified, the residual inertia characters have orders
$134$ at $r=66$ and $67$ at $r=200$, both nontrivial.
An inertia scalar therefore differs from one by an $O$-unit,
and $D^{G_{K_p}}=0$.
At a finite place away from $p$, the cyclotomic, Teichmuller
and $\kappa$ factors are unramified.
The finite inertia action has order dividing four, prime to $p$;
its invariants are either zero or the whole divisible module $D$.
Thus $D^J/\varpi D^J=0$ for every local group in (P16Q.1).

The multiplication-by-$\varpi$ long exact sequence gives
local injectivity $H^1(J,D[\varpi])\to H^1(J,D)$.
A strict global class killed by $\varpi$ lifts to a residual
class, which remains strict by local injectivity and is zero
by Step~6.
Every continuous cocycle from compact $G_K$ to discrete $E/O$
has finite image killed by a power of $\varpi$.
A nonzero strict class would yield a nonzero one killed by
$\varpi$, a contradiction.  This proves (P16Q.2) without
changing the original coefficient ring.
\end{proof}

\paragraph{Actual labels, count and geometric scope.}
For $m=0,\kappa=1$, the labels are exactly
$\phi|_{G_K}$ and $\phi^{-1}|_{G_K}$, with trivial cubic part.
They are distinct because $K$ is disjoint from the conductor-sixteen
quartic field, and they occur in every prescribed host of (FG.3)
with $h\geq2$, since that host contains $\mathbb Q(\zeta_{16})$.
They are distinct from the order-one and order-two labels and
from all labels with a nontrivial cubic restriction.
The new weights are $e=66+134n$, $n\geq0$.
For positive integer $T$, the exact total count over $1\leq e<T$ is
$2\lfloor(T+67)/134\rfloor$.
No half-residue strict assertion for these pure quartic labels
is inferred from this theorem.
The finite-cover/rank-one local-system dictionary is the retained
SGA~1, Expos\'e~V, \S\S4 and 7 dictionary; Kummer torsors retain
the audited Stacks tag~03PK scope.
The residual covers used in the proof do not enlarge the list
of original finite labels.  Positive derived indices and the
full local comparison cost remain separate; actual window
transport is not asserted here.

\paragraph{Evidence and applied-source scope.}
The complete exact verifier is
\path{verify_pure_conductor16_quartic.py}, with output
\path{pure_conductor16_quartic_certificate.json}.
It checks both weighted sums and integral floor sums, the four
annihilator units, the exact field automorphisms, the two global
unit relations and all four nonlinear local projectors.
The original AC-G147 unit files and class proof are copied unchanged.
The applied field-character proof (G147C.1)--(G147C.10) and its
class theorem are retained at their recorded whole-proof scope;
no repeated analytic interval execution is claimed.
Sharifi's complete Definition~4.5.1 statement, Lemma~4.5.2 proof
and Theorem~4.5.6 proof were reread from the retained primary
chapter, with exact passages copied in this evidence folder.
The earlier Gauss lemmas inside that published theorem, and
the class norm/extension and equivariant unit theorems, remain
published or previously audited inputs at their precise scopes.
No mixed realization, all-prime estimate, H04.4 or Lean execution
is claimed.
